\section{Определение ряда из комплексных чисел. Определение сходящегося
ряда. Связь со сходимостью рядов действительной и мнимой
частей.}\label{ux43eux43fux440ux435ux434ux435ux43bux435ux43dux438ux435-ux440ux44fux434ux430-ux438ux437-ux43aux43eux43cux43fux43bux435ux43aux441ux43dux44bux445-ux447ux438ux441ux435ux43b.-ux43eux43fux440ux435ux434ux435ux43bux435ux43dux438ux435-ux441ux445ux43eux434ux44fux449ux435ux433ux43eux441ux44f-ux440ux44fux434ux430.-ux441ux432ux44fux437ux44c-ux441ux43e-ux441ux445ux43eux434ux438ux43cux43eux441ux442ux44cux44e-ux440ux44fux434ux43eux432-ux434ux435ux439ux441ux442ux432ux438ux442ux435ux43bux44cux43dux43eux439-ux438-ux43cux43dux438ux43cux43eux439-ux447ux430ux441ux442ux435ux439.}

\textbf{Определение:} Пусть \(\{z_n\}\), \(z_n \in \mathbb{C}\).
Числовой ряд:
\[\sum_{n=1}^{\infty} z_n = z_1 + z_2 + \dots, \qquad z_n \text{ — общий член.}\]

\textbf{Определение:} Частичная сумма:
\[S_n = \sum_{k=1}^{n} z_k, \qquad \{S_n\} \text{ — последовательность частичных сумм.}\]

\textbf{Определение:} Ряд \(\sum z_n\) сходится
\(\;\Leftrightarrow\; \exists\, S \in \mathbb{C}:\)
\[\forall \varepsilon > 0 \;\; \exists n_\varepsilon \;\; \forall n \geq n_\varepsilon : \; |S_n - S| < \varepsilon.\]
\(S = \lim\limits_{n\to\infty} S_n = \sum\limits_{n=1}^{\infty} z_n\)
--- сумма ряда. Иначе ряд расходится.

\textbf{Теорема (связь с Re/Im):} \(z_n = a_n + i b_n\),
\(\;a_n, b_n \in \mathbb{R}\):
\[\sum_{n=1}^{\infty} z_n \text{ сходится} \;\;\Longleftrightarrow\;\; \exists \sum_{n=1}^{\infty} a_n \;\wedge\; \exists \sum_{n=1}^{\infty} b_n,\]
\[\sum_{n=1}^{\infty} z_n = \sum_{n=1}^{\infty} a_n + i \sum_{n=1}^{\infty} b_n.\]
